\chapter{Contribution}

\section{Motivation}

\begin{remark}
Railway alignment has traditionally been approached as a geometric
design problem, supplemented by engineering rules and dynamic
considerations.
\end{remark}

\begin{remark}
However, these aspects are often treated separately, leading to
fragmented models and limited theoretical coherence.
\end{remark}

---

\section{Central Hypothesis}

\begin{theorem}[Curvature as System Variable]
A railway alignment can be represented as a curvature-driven system in
which
\[
\kappa(s)
\]
acts as the primary state variable.
\end{theorem}

\begin{remark}
All geometric, kinematic, and dynamic properties of the alignment follow
from $\kappa(s)$ and its derivatives.
\end{remark}

---

\section{System Interpretation of Alignment}

\begin{proposition}
The behaviour of a railway alignment can be described by the mapping
\[
\kappa(s) \;\longrightarrow\; \text{forces} \;\longrightarrow\;
\text{system response} \;\longrightarrow\; \text{performance}.
\]
\end{proposition}

\begin{remark}
This establishes alignment as a system rather than a static geometric
object.
\end{remark}

---

\section{Unification of Transition Curves}

\begin{theorem}[Unified Transition Representation]
All transition curves can be expressed as curvature functions
\[
\kappa(s)
\]
within a normalized domain.
\end{theorem}

\begin{remark}
This allows comparison of historically distinct transition concepts
within a single mathematical framework.
\end{remark}

---

\section{Transition Curves as Control Functions}

\begin{theorem}[Control Interpretation]
Transition curves act as control functions governing the evolution of
curvature and thus the dynamic behaviour of the system.
\end{theorem}

\begin{remark}
This interpretation replaces the traditional view of transition curves
as purely geometric connectors.
\end{remark}

---

\section{Structured Alignment Model}

\begin{proposition}
An alignment can be represented as a structured sequence of elements:
\begin{itemize}
\item fixed curvature elements,
\item transition elements,
\item auxiliary alignment layers such as cant.
\end{itemize}
\end{proposition}

\begin{remark}
This structure separates representation, geometry, and system behaviour.
\end{remark}

---

\section{Coupling with Engineering Variables}

\begin{remark}
Alignment design is governed by coupled variables including:
\begin{itemize}
\item curvature,
\item cant,
\item speed,
\item and their rates of change.
\end{itemize}
\end{remark}

\begin{proposition}
Dynamic constraints can be expressed as conditions on
\[
\kappa(s), \quad \kappa'(s), \quad \kappa''(s).
\]
\end{proposition}

---

\section{Numerical and Optimization Perspective}

\begin{remark}
The curvature-based representation allows direct application of
numerical integration and optimization methods.
\end{remark}

\begin{proposition}
Alignment design can be formulated as a parameterized optimization
problem in curvature space.
\end{proposition}

---

\section{Integration with Data Models}

\begin{remark}
Modern railway data models represent alignment as a multi-layer
structure consisting of horizontal, vertical, and cant components.
\end{remark}

\begin{proposition}
The structured alignment model proposed here is compatible with such
multi-layer representations.
\end{proposition}

---

\section{Main Contributions}

\begin{summary}
The contributions of this work can be summarized as follows:

\begin{itemize}
\item A curvature-based system interpretation of alignment.
\item A unified representation of transition curves.
\item The interpretation of transition curves as control functions.
\item A structured alignment model separating geometry and behaviour.
\item Integration of geometry, dynamics, and engineering constraints.
\item A foundation for alignment-based information modelling (AIM).
\end{itemize}

The numerical comparison of transition laws shows that identical
boundary conditions may still lead to substantially different internal
curvature evolution and therefore different dynamic behaviour.
\end{summary}

---

\section{Scientific Positioning}

\begin{remark}
The presented approach does not introduce a new transition curve, but a
new conceptual framework for understanding alignment.
\end{remark}

\begin{remark}
It connects previously separated domains:
\begin{itemize}
\item geometry,
\item dynamics,
\item engineering design,
\item and data modelling.
\end{itemize}
\end{remark}

---

\section{Outlook}

\begin{remark}
The framework provides a basis for future work in:
\begin{itemize}
\item dynamic simulation,
\item alignment optimization,
\item network-level modelling,
\item and BIM integration.
\end{itemize}
\end{remark}
